\documentclass[11pt]{article}
\usepackage{amsmath}
\usepackage{amssymb}
\usepackage[margin=2cm]{geometry}
\usepackage{longtable}
\usepackage{booktabs}
\usepackage{pdflscape}
\usepackage{breqn}
\title{Certificate for the q-integer [5]\_q}
\author{qreals certificate}
\date{\today}
\begin{document}
\maketitle

\noindent Input: the integer n = 5.

\section*{(a) Continued fraction and even-length MGO form}
Regular continued fraction: [5]. \par
Even-length MGO form: [5]. \par
The even-length form evaluates to the convergent $ 5 $.\par

\section*{(b) MGO formula folded step by step}
Odd positions carry $[a]_q$ with $q^a$ above; even positions carry
$[a]_{q^{-1}}$ with $q^{-a}$ above (MGO eqn.~1.1).

\noindent\textbf{[5]\_q is the Gauss q-integer 1 + q + ... + q\textasciicircum{}(n-1):}
{\footnotesize%
\begin{dmath*}
q^{4} + q^{3} + q^{2} + q + 1
\end{dmath*}
}

\noindent\textbf{Result.} [5]\_q:
{\footnotesize%
\begin{dmath*}
q^{4} + q^{3} + q^{2} + q + 1
\end{dmath*}
}

\noindent Taylor coefficients:
{\small%
\begin{dmath*}
1 + q + q^{2} + q^{3} + q^{4} + O(q^{12})
\end{dmath*}
}

\section*{(c) Independent cross-checks}
\begin{itemize}
\item \textbf{[pass]} [5]\_q at q=1 is 5, the ordinary value 5
\item \textbf{[pass]} the Taylor expansion of the exact rational function and the truncated series agree on q\textasciicircum{}0..q\textasciicircum{}11: [1, 1, 1, 1, 1, 0, 0, 0, 0, 0, 0, 0]
\end{itemize}
Summary: verified: q=1 matches 5, exact = truncated to 12.

\section*{References}
S. Morier-Genoud and V. Ovsienko, q-deformed rationals and q-continued
fractions, Forum Math. Sigma \textbf{8} (2020), e13. Per-function theorem and
check mapping: \texttt{docs/CORRECTNESS.md}.

\medskip
\noindent This is a human-auditable derivation, not a formal machine proof;
every line above is checkable by hand.
\end{document}
